\documentclass{article}

\IfFileExists{neurips_2026.sty}{%
  \usepackage[final]{neurips_2026}%
}{%
  \usepackage[margin=1in]{geometry}%
  \PackageWarningNoLine{main}{neurips_2026.sty not found - using fallback geometry. Page count will NOT match the official template.}%
}

\makeatletter
\renewcommand{\@noticestring}{NewInML Workshop at the 40th Conference on Neural Information Processing Systems (NeurIPS 2026), Paris, France, 2026.}
\makeatother

\usepackage[utf8]{inputenc}
\usepackage[T1]{fontenc}
\usepackage{hyperref}
\usepackage{url}
\usepackage{booktabs}
\usepackage{amsfonts}
\usepackage{nicefrac}
\usepackage{microtype}
\usepackage{xcolor}
\usepackage{graphicx}
\usepackage{amsmath}

\newcommand{\justification}[1]{#1}

\graphicspath{{../results/figures/}{figures/}}

\title{Supplementary Material for\\ Metric-as-Feature Leakage: A Feature-Side\\
       Failure Mode in Retriever Routing}

\author{
  Soumy Sawant\\
  Manipal University Jaipur\\
  Jaipur, India\\
  \texttt{soumy.250205119@muj.manipal.edu}
  \And
  Chirag Joshi\\
  Manipal University Jaipur\\
  Jaipur, India\\
  \texttt{joshi.chirag@jaipur.manipal.edu}
}

\begin{document}
\maketitle
\renewcommand{\thesection}{S\arabic{section}}
Section, table and figure numbers without an S refer to the main paper.

\section{Positive control: a genuine gain beside a leak}
\label{sup:positive}
Our data has no exploitable honest signal, so we build a synthetic task that does: three of twenty Gaussian features set each arm's
utility, and rewards are reciprocal-rank-like. With no signal, gradient boosting matches the
majority rate ($-0.4$ points) and \texttt{leakcheck} raises nothing. With signal it gains $+14.0$
and $+17.4$ points at two strengths, with no verdict. Appending each arm's reward with independent
additive noise of $\pm$0.15 adds a further 26--34 points, again with no verdict. Suspicions, however, fall on
the genuine signal features as well as on the leaked ones. A gain is therefore not evidence of
leakage, and neither is a suspicion: both are resolved by ablation, as \S6 of the paper argues.


\section{Supporting checks for Section~5 of the paper}
\label{app:checks}

\paragraph{Why the label is hard.} The reward depends on where a gold passage landed, and which
passage counts as gold depends on the candidate draw: 34.3\% of the 1{,}989 prepared NQ queries
contain at least one answer-bearing distractor (13.4 on average among those that do; 4.59 over
all), as do 46.1\% of TriviaQA and 27.0\% of HotpotQA queries, so the gold rank is partly
arbitrary. Gold is the presence of the answer string in a 100K-passage Wikipedia sample, which
1{,}157 of 1{,}989 NQ, 1{,}539 of 1{,}894 TriviaQA and 916 of 1{,}995 HotpotQA queries have. Systematic BM25--Dense differences exist (rare terms
versus paraphrase) but are smaller than this variance. The oracle over more arms is also a maximum over more noise, so more arms of the same
kind would raise the headroom without necessarily raising what a selector can capture; a more
diverse pool remains the proper test (Limitations of the paper).


\paragraph{Removing the redundant arm.} Hybrid loses to Dense everywhere (\S3 of the paper), so the
failure might come from a non-competitive action. Restricting to
BM25 and Dense moves lift from $-3.3/-3.1/-6.3$ points to $-2.0/-2.8/-5.0$: still negative. This does not show that the arm caused the failure, because dropping Hybrid also
raises the majority baseline (63.4$\to$69.6, 61.7$\to$69.0, 56.4$\to$63.0) and changes the informative set,
and a two-class problem at 70\% majority sits closer to its baseline than a three-class one at
63\%. Detection is also harder here: on the test splits a router would need to capture 85\% (NQ), 51\%
(TriviaQA) or 78\% (HotpotQA) of the two-arm headroom to be detectable (one-sided, $\alpha{=}0.05$,
80\% power).


\paragraph{The structural findings on RouterBench.} The two diagnostics that \S5 of the paper
applies to our data give, on RouterBench, 69.9\% of oracle headroom from queries won by
$\le 1\%$ margins (ours: 16--22\%) and 40\% headroom inflation under noise at 10\% of the reward's
spread (ours: 16--29\%). Same sign, larger magnitude, over 11 arms rather than three, as expected when a maximum is taken
over more arms.

\paragraph{Per-dataset null cells.} Observed vs.\ shuffled-label null ($n$, majority): NQ 60.1
vs.\ $59.1{\pm}0.6$, $p{=}0.066$ (1{,}021, 63.4); TriviaQA 58.6 vs.\ $58.2{\pm}0.5$, $p{=}0.181$
(1{,}307, 61.7); HotpotQA 50.1 vs.\ $50.8{\pm}1.0$, $p{=}0.740$ (847, 56.4). Random-feature null:
0.090, 0.187, 0.820. On HotpotQA the observed point is below its own null. NQ has the largest per-dataset effect and is
still not significant.

\paragraph{Baseline details for Table~3 of the paper.} The score-gap heuristic was repaired before
comparison: it swept a fixed threshold range outside the feature's support, and sweeping the
training split's quantiles instead raises it from 0.1029 to 0.1657 on NQ. Best fixed is chosen in
hindsight on the test split; chosen on training data it would be Hybrid on HotpotQA, against which
every method ``gains'' 12--35\%, all by always choosing Dense. Convex fusion, which combines rather
than chooses, gains on TriviaQA ($+7.0\%$ capture) and loses on NQ ($-2.4\%$) and HotpotQA
($-45.2\%$). Paired-bootstrap 95\% intervals on capture are wide: score-gap on NQ is $+0.4\%$
[0.0, 1.5], and the largest upper limit among the supervised methods is 45.6\% (ridge, HotpotQA),
so point estimates are not ceilings. RL capture per dataset is $-10.8$, $-0.0$ and $-9.1\%$
(single seeds $-49\%$ to $+9\%$). Two arms (BM25 and Dense) would need 85\%, 51\% and 78\% capture
to be detectable (one-sided, 80\% power).

\paragraph{Policy-gradient settings and additional negatives.}
REINFORCE~\cite{williams1992reinforce}, PPO~\cite{schulman2017ppo} and
GRPO~\cite{shao2024deepseekmath} span 0.0139 MRR on NQ, all below best-fixed---a separate
five-seed sweep at 300 epochs, reporting final-epoch weights, rather than the fifteen-seed
best-validation run of Table~3 of the paper, so the two are not directly comparable.
The policy of Table~3 of the paper is a three-layer MLP (hidden width 128) trained with REINFORCE
for 150 epochs at learning rate $3{\times}10^{-3}$, entropy bonus 0.01 and a per-query advantage
(each arm's reward minus the query's mean over arms), on the full reward table of the training
split; the checkpoint is chosen on the validation split. A cost-aware cascade finds no Pareto improvement; gold-type
stratification changes nothing. Per-retriever reward regression, scored as MRR@20, is negative on 11 of 15
configurations; the exception is HotpotQA, where ridge beats best-fixed by $+0.021$ and $+0.024$
($p{=}0.036$, $p{=}0.029$). Those are the only nominal $p<0.05$ results among the fifteen, they
sit on the smallest split, and they are ridge variants differing only in whether embeddings are
included---one finding rather than two. A cross-score feature variant reproduces both and adds a
third marginal hit ($p{=}0.046$). None survives Bonferroni ($\alpha{=}0.0033$) nor reaches the
72\% capture that detection on the HotpotQA test split would require (one-sided, 80\% power), so we
do not claim them.



\section{Detector benchmark: which check catches which leak}
\label{sup:bench}
Each synthetic problem has 20 honest features, three of which carry genuine signal about the
winning arm, and reciprocal-rank-like rewards; one planted column (or block) is added per scenario,
built from arm 0's reward. A detector flags a column if it fires against any arm. Ten seeds per
scenario, $n{=}2000$. The leaks copy a full reward and are therefore strong; rates near 1 reflect
that, not performance on subtler leaks.

\begin{table}[h]
\centering
\small
\caption{Rate at which each detector flags the planted column. ``none'' and ``confound'' are not
leaks, so their rates are false flags. Equality: exact agreement on more than 25\% of non-zero
rows. AUC scan: bioLeak-style $|\mathrm{AUC}-0.5|\times 2 \ge 0.90$. Importance: held-out
permutation importance of a gradient-boosting router $>0.01$ accuracy (not SHAP). v1: any
\texttt{leakcheck} 0.1 finding; v2: any finding of the opt-in decodability test.}
\label{tab:bench}
\begin{tabular}{lccccccc}
\toprule
Scenario & Equality & Pearson & AUC scan & Importance & v1 verdict & v1 any & v2 any \\
\midrule
none (not a leak)          & 0 & 0   & 0   & 0.2 & 0 & 0   & 0 \\
exact copy                 & 1 & 1   & 0.2 & 1   & 1 & 1   & 1 \\
truncated (MRR@5-like)     & 1 & 1   & 0   & 1   & 1 & 1   & 1 \\
monotone (rank)            & 0 & 1   & 0.2 & 1   & 0 & 1   & 1 \\
value-keyed $\pm$0.15      & 0 & 1   & 0   & 1   & 0 & 1   & 1 \\
order-scrambled recoding   & 0 & 0.2 & 0   & 1   & 0 & 0.8 & 1 \\
independent noise $\pm$0.15 & 0 & 1  & 0   & 1   & 0 & 1   & 1 \\
independent noise $\pm$0.50 & 0 & 1  & 0   & 1   & 0 & 1   & 1 \\
group confound (not a leak) & 0 & 0  & 0   & 1   & 0 & 0.9 & 1$^\dagger$ \\
split over two columns     & 0 & 1   & 0   & 1   & 0 & 1   & 1 \\
PCA-mixed over six columns & 0 & 1   & 0   & 1   & 0 & 1   & 1 \\
\midrule
genuine-signal columns     & 0 & 0.67 & 0  & 0.85 & 0 & 0.74 & 0.91 \\
pure-noise columns         & 0 & 0    & 0  & 0.026 & 0 & 0   & 0 \\
\bottomrule
\end{tabular}

{\footnotesize $^\dagger$ labelled \texttt{GROUP CONFOUND} in 10 of 10 seeds when the group id is supplied.}
\end{table}

Verdicts never fire on a non-leak, but they catch only exact and truncated copies. Every detector
that catches perturbed leaks also flags genuinely predictive features, because a predictive
feature is what a leak looks like from the data; flags are triage, resolved by ablation. The
decodability test adds three things: it catches the order-scrambled recoding in 10 of 10 seeds
(against 8 of 10 for v1 and 2 of 10 for Pearson), it names the group confound as one, and it raises
no flags on pure-noise columns. On our own pooled data it flags the leak slots 22--23 under the
value-keyed recoding at $\pm$50\% and $\pm$100\%, where v1 flags nothing, and on the honest matrix it
flags the label-free overlap features 26--28, which carry real signal. None of its real-data flags
reads as a group confound.

\bibliographystyle{plain}
\bibliography{references}

\end{document}
